On the Fock space of $\nu \geq 2$ boundary modes the bivectors of $\so(2\nu)$ split as hops $a_i^{\dagger}a_j$ (with the Cartan elements $a_i^{\dagger}a_i - \tfrac12$), pair annihilations $a_ia_j$ and pair creations $a_i^{\dagger}a_j^{\dagger}$. On the vacuum every hop with $i \neq j$ and every annihilation vanishes and every Cartan element is a scalar, so hops and annihilations are the stabiliser of the vacuum ray; the creations span the tangent of the orbit of Witt vacua, of dimension $\nu(\nu-1)/2$. A hop rotation fixes the vacuum, a pairing rotation moves it, and the necessity witness of the singularity paper's envelope corollary, $\ket{00} \mapsto \tfrac35\ket{00} + \tfrac45\ket{11}$, is the pairing rotation on one pair at $(\cos\theta, \sin\theta) = (\tfrac35, \tfrac45)$. \emph{Provenance:} \cite[Prop.~13.4, Cor.~11.6]{ButlerT5} (\claimref{singularity.envelope-endpoint-why-boundary-mode-locality}, \claimref{singularity.l-stabiliser-tangent}); the orbit is the spinor variety of \cite{Cartan1938, Chevalley1954}, imported there, and the exponential form of a rotated vacuum is Thouless's \cite{Thouless1960}, imported by the laboratory paper \cite[Fact~2.10]{ButlerT16} (\claimref{one-relation.imported-cartans-purity-theorem-thouless-form}).

This restates \claimref{pairing.imported-stabiliser-tangent-external}.
